\documentclass[12pt]{letter}
\usepackage[T1]{fontenc}
\usepackage[utf8]{inputenc}
\usepackage[paper=a4paper, left=2.5cm, right=2.5cm, top=2.5cm, bottom=2.5cm]{geometry}
\usepackage{times}
\usepackage{hyperref}
\usepackage{parskip}

\signature{Roberto S\'{a}nchez-Reolid}

\address{Instituto de Investigaci\'{o}n en Inform\'{a}tica de Albacete (I3A)\\
Universidad de Castilla-La Mancha\\
Campus Universitario s/n\\
02071 Albacete, Spain\\
\href{mailto:roberto.sanchez@uclm.es}{roberto.sanchez@uclm.es}}

\begin{document}

\begin{letter}{The Editor-in-Chief\\
Biomedical Signal Processing and Control\\
Elsevier}

\opening{Dear Editor-in-Chief,}

We submit our manuscript entitled \textbf{``A Comparative Study of Emerging Architectures for EDA-based Arousal Classification''} for consideration for publication in \textit{Biomedical Signal Processing and Control}.

Electrodermal Activity (EDA) is a non-invasive biomarker of sympathetic arousal increasingly used in wearable affective computing. While Transformer-based architectures outperform convolutional baselines for subject-independent arousal classification, their computational cost limits deployment on latency-sensitive platforms. This work presents a systematic comparison of eight emerging architectures spanning five paradigms---patch-based attention, sparse attention, frequency-domain modelling, state space models, and modernised convolutions---for binary arousal classification from EDA signals under strict Leave-One-Subject-Out (LOSO) validation with 147 participants. Beyond predictive accuracy, we introduce five complementary computational efficiency metrics---parameter count, FLOPs, inference time, memory footprint, and training time---as primary evaluation dimensions, enabling Pareto-frontier characterisation.

Mamba, a selective state space model, achieves F1 = 0.858 at 1.5\,ms inference time, statistically indistinguishable from PatchTST accuracy (F1 = 0.863) with 3.6$\times$ lower latency and 66\% fewer parameters. These results directly inform architecture selection under deployment constraints in wearable, edge, and server scenarios.

We believe this manuscript is well-suited for BSPC for the following reasons:
\begin{enumerate}
  \item \textbf{Methodological contribution:} The paper provides the first systematic comparison of Transformer, state space model (SSM), and modernised convolutional architectures for EDA-based arousal classification under strict LOSO validation. By elevating computational efficiency to a primary evaluation dimension alongside classification performance, the work addresses a practical gap between laboratory accuracy and real-world deployability.
  \item \textbf{Scope alignment:} The manuscript directly addresses BSPC's scope in wearable health monitoring systems, machine learning for biomedical signal analysis, real-time embedded implementations, and data-driven health analytics. The explicit characterisation of latency constraints across wearable/edge/server tiers speaks to the journal's emphasis on translational relevance.
  \item \textbf{Technical novelty:} The Pareto-frontier characterisation with five complementary efficiency metrics (parameter count, FLOPs, inference time, memory, and training time) is, to our knowledge, the most comprehensive efficiency evaluation of deep architectures for physiological signal classification published to date.
  \item \textbf{Reproducibility:} All experiments use a publicly available dataset with 147 participants. The complete codebase, preprocessing pipeline, and hyperparameter configurations have been made available to ensure full reproducibility.
\end{enumerate}

We confirm that:
\begin{itemize}
  \item This manuscript has not been published previously and is not under consideration elsewhere.
  \item All authors have read and approved the final version of the manuscript.
  \item The authors declare no competing financial interests.
  \item A CRediT author statement and funding declaration are included in the manuscript.
  \item A declaration of generative AI use in manuscript preparation is included.
\end{itemize}

We appreciate your time and consideration and look forward to your decision.

\closing{Yours sincerely,}

\end{letter}
\end{document}
